% !PRINTED 315
\weberpage{315}
%<<CONT
Feldern giebt also der Verlust an Zeichenwechseln die richtige Anzahl von
Wurzeln. In dem Felde 2 dagegen ist die Entscheidung zwischen zwei und
keiner Wurzel nicht getroffen.

Nehmen wir nun das \gesperrt{Jacobi}'sche Kriterium (\S.\,102). Danach
müssen wir, um die Anzahl der zwischen $\alpha$ und $\beta$ gelegenen
Wurzeln abzuschätzen, die Gleichung (1) erst transformiren auf
\[
(\alpha+\beta y)^2+a(\alpha+\beta y)(1+y)+b(1+y)^2=0
\]
oder
\[
\begin{aligned}
&(\beta^2+a\beta+b)y^2+(2\alpha\beta+a(\alpha+\beta)+2b)y\\
&\quad{}+(\alpha^2+a\alpha+b)=0,
\end{aligned}
\]
und es ist nun die Anzahl der Zeichenwechsel in den drei Functionen
\[
\begin{aligned}
&\beta^2+a\beta+b\\
&2\alpha\beta+a(\alpha+\beta)+2b\\
&\alpha^2+a\alpha+b
\end{aligned}
\]
abzuzählen. Der erste und der dritte dieser Ausdrücke stellen, gleich Null
gesetzt, die beiden vorhin betrachteten Parabeltangenten dar. Der mittlere,
$2\alpha\beta+a(\alpha+\beta)+2b$, verschwindet für $a=-2\alpha$,
$b=\alpha^2$ und für $a=-2\beta$, $b=\beta^2$ stellt also die
Verbindungslinie der beiden Punkte $\alpha$, $\beta$ dar, und zwar so,
dass er auf der Seite, die den Schnittpunkt der beiden Tangenten enthält,
negativ wird [in diesem Schnittpunkte selbst ist er gleich
$-(\beta-\alpha)^2$].

% SOURCE FIGURE CROP: Bd. I registered scan 341 (printed p. 315), 600 dpi,
% crop x=100,y=3000,w=700,h=620; source-faithful raster, Fig. 20.
\weberfigureimage{weber-b1-fig20.png}{0.34}
Wir erhalten hier fünf Felder mit keinem, einem oder zwei
Zeichenwechseln, wie die Fig. 20 angiebt.

Man sieht hieraus, dass diese Figur die Unentschiedenheit auf einen viel
kleineren Raum beschränkt als die vorige. Von diesem Gesichtspunkte aus
erscheint es also nicht gerechtfertigt, wenn \gesperrt{Jacobi} dem
\gesperrt{Budan-Fourier}'schen Kriterium vor dem seinigen so entschieden
den Vorzug giebt\footnote{Vgl. F.\ \gesperrt{Klein}, Geometrisches zur
Abzählung der reellen Wurzeln einer algebraischen Gleichung in
\gesperrt{Dyck}'s Catalog mathematischer Modelle (München 1892).}.
